\section{Introdução}

O câncer do colo do útero permanece entre as principais causas evitáveis de
morbidade e mortalidade por câncer em mulheres. Estimativas globais para 2020
indicam aproximadamente 604 mil novos casos e 342 mil mortes, com maior carga
em regiões que enfrentam restrições de acesso à vacinação, ao rastreamento e ao
tratamento oportuno \cite{singh2023global}. A citologia cervical continua sendo
um componente relevante de programas de rastreamento, mas a leitura manual de
lâminas é trabalhosa, depende de especialistas e envolve alterações
morfológicas sutis.

Redes neurais profundas têm apresentado resultados promissores na análise de
imagens citológicas \cite{jiang2023systematic,fang2024systematic}. Entretanto, a
comparação entre métodos é dificultada por diferenças de base, preparação da
lâmina, tarefa, granularidade dos rótulos e protocolo de validação
\cite{fang2024systematic}. Um risco central ocorre quando células recortadas da
mesma imagem aparecem em treino e avaliação: elas compartilham coloração,
aquisição, fundo e artefatos, podendo superestimar a generalização para imagens
não vistas.
Por isso, diretrizes de relato para modelos preditivos e inteligência
artificial em imagens médicas recomendam explicitar fonte dos dados, unidade de
análise, nível de separação entre conjuntos e uso pretendido do sistema
\cite{collins2024tripod,tejani2024claim}.

Este artigo estuda a classificação binária \emph{normal versus anormal} de
células da coleção CRIC Cervix. ASC-US, LSIL, ASC-H, HSIL e SCC compõem a
classe anormal, enquanto NILM compõe a classe normal. A contribuição principal é
avaliar a ConvNeXt-Tiny \cite{liu2022convnext} com validação cruzada agrupada
por imagem-fonte, reportando média, desvio-padrão e variação fold a fold. Para
contextualizar o desempenho, inclui-se a ResNet-50 \cite{he2016resnet} como
comparador complementar treinado nos mesmos folds.
